\documentclass[../../main.tex]{subfiles}
\begin{document}

\begin{flushright}
  \footnotesize\textit{【自動文字起こし・要確認】}
\end{flushright}

% Source: handwritten.pdf 第1問．
{\bf [解答]}
$f(x)=px+q-\log x$とおく．$p\leqq0$の時，$f$は連続で
\begin{align}
  f(x)\longrightarrow+\infty\quad(x\to+0),\qquad
  f(x)\longrightarrow-\infty\quad(x\to+\infty)
\end{align}
から，$f(x)=0$なる$x$が存在し不適．よって$p>0$としてよい．
\begin{align}f'(x)=p-\dfrac{1}{x}
\end{align}
より下表である．
\begin{table}[htb]
  \centering
  \begin{tabular}{c|ccccc}
    $x$ & $0$ & & $\dfrac{1}{p}$ & & $+\infty$ \\
    \hline
    $f'(x)$ & & $-$ & $0$ & $+$ & \\
    \hline
    $f(x)$ & $+\infty$ & $\searrow$ & $f(1/p)$ & $\nearrow$ & $+\infty$
  \end{tabular}
  \caption{$p>0$のときの$f(x)$の増減表}
\end{table}
よって，共有点を持たない条件は
\begin{align}
  f(x)>0\ \Longleftrightarrow\ f\left(\dfrac{1}{p}\right)>0
  \ \Longleftrightarrow\ 1+q+\log p>0\qquad(p>0).
\end{align}

\end{document}
